\section{Results}
\label{sec:results}

\subsection{Convergence on geodesics}
\label{sec:convergence}

Every fixed-step method reaches its theoretical order on geodesics, in both formulations
(\fig{convergence}, \tab{orders}). On the strong-field orbit $(7.5, 0.5)$ the measured orders
are within $\pm0.05$ of theory, on the deflection with $b = 6$ within $\pm0.08$. GL1 and Tao2
are almost indistinguishable in (b), and GL1 is about 35 times less accurate in (a) than in (b)
at the same step. The adaptive pairs show tolerance proportionality: the global error falls
like $\mathrm{tol}^{0.85\text{--}1.04}$ for DP5 and DOP853, on the same curves as SciPy's
(\fig{tolerance}). On $(20, 0.5)$ RK4 stays pre-asymptotic over the whole usable range
(slopes 4.5 to 4.07 before round-off), which is why the orders are measured on the
strong-field orbit.

\begin{figure}[t]
\centering
\includegraphics[width=0.95\textwidth]{F04_convergence.pdf}
\caption{\label{fig:convergence}Convergence of the fixed-step methods: error against step size
for the deflection $b = 6$ (error in $\Delta\phi(r_\mathrm{far})$) and one radial period of
the orbit $(7.5, 0.5)$ (error in $\phi$), in both formulations, with reference slopes 2, 4
and 6 (F4).}
\end{figure}

\begin{table}[t]
\caption{\label{tab:orders}Measured order of convergence (T3), least-squares slope over the
asymptotic range with its standard error.}
\centering\small
% Generated by scripts/make_tables.py from results/summary/. Do not edit.
\begin{tabular}{lccccc}
\toprule
Method & Theory & Deflection (a) & Deflection (b) & Orbit (a) & Orbit (b) \\
\midrule
RK4 & 4 & $4.08 \pm 0.02$ & $3.96 \pm 0.01$ & $3.98 \pm 0.01$ & $4.01 \pm 0.00$ \\
GL1 & 2 & $2.00 \pm 0.00$ & $2.00 \pm 0.00$ & $2.00 \pm 0.00$ & $2.00 \pm 0.00$ \\
GL2 & 4 & $4.00 \pm 0.00$ & $3.98 \pm 0.01$ & $3.98 \pm 0.01$ & $4.00 \pm 0.00$ \\
GL3 & 6 & $6.00 \pm 0.00$ & $6.06 \pm 0.04$ & $5.97 \pm 0.01$ & $6.05 \pm 0.01$ \\
Tao2 & 2 & -- & $2.00 \pm 0.00$ & -- & $2.00 \pm 0.00$ \\
Tao4 & 4 & -- & $4.01 \pm 0.00$ & -- & $4.00 \pm 0.00$ \\
\bottomrule
\end{tabular}

\end{table}

\begin{figure}[t]
\centering
\includegraphics[width=0.6\textwidth]{F05_tolerance.pdf}
\caption{\label{fig:tolerance}Tolerance proportionality of the adaptive pairs: global error
against $\mathrm{rtol} = \mathrm{atol}$, ours and SciPy's (F5).}
\end{figure}

\subsection{Short integrations: cost per accuracy (RQ3)}
\label{sec:short}

\fig{work-precision} and \tab{cost} compare the cost of reaching a given error on three short
problems: the deflection $b = 6$, three periods of $(20, 0.5)$ and two periods of the same
orbit inclined by $60^\circ$. \textbf{DOP853 is the cheapest method on every problem and at
every accuracy.} To reach $10^{-10}$ in formulation (b) it needs $1.4$--$3.2\times10^3$
evaluations, against $1.6$--$3.6\times10^4$ for GL3, the best fixed-step method, and
$3.5\times10^4$--$1.7\times10^5$ for RK4. Tao4 costs 1.2--10 times as much as RK4 for the same
error, and GL1 and Tao2 never reach $10^{-10}$ in the sweep. Formulation (b) is the cheaper
one for most pairs of method and problem, but not all: on the inclined orbit DOP853 and GL3
reach $10^{-10}$ more cheaply in (a).

\begin{figure}[t]
\centering
\includegraphics[width=0.9\textwidth]{F06_work_precision_b.pdf}
\caption{\label{fig:work-precision}Work-precision in formulation (b): error against
vector-field evaluations (left) and wall time (right) for the three short problems (F6).
Each curve is the Pareto front of a sweep over 12 step sizes or 11 tolerances.}
\end{figure}

\begin{table*}[t]
\caption{\label{tab:cost}Evaluations needed to reach the error $10^{-6}$ or $10^{-10}$ (T5),
interpolated along each method's work-precision front, in formulation (b) with (a) in
parentheses. A dash means the target was not reached in the sweep.}
\centering\small
\resizebox{\textwidth}{!}{% Generated by scripts/make_tables.py from results/summary/. Do not edit.
\begin{tabular}{lrrrrrr}
\toprule
Method & Deflection, $10^{-6}$ & $10^{-10}$ & Orbit, $10^{-6}$ & $10^{-10}$ & Inclined, $10^{-6}$ & $10^{-10}$ \\
\midrule
RK4 & $1.7\times10^{4}$ ($1.7\times10^{4}$) & $1.7\times10^{5}$ ($1.6\times10^{5}$) & $4.8\times10^{3}$ ($8.3\times10^{3}$) & $3.5\times10^{4}$ ($8.4\times10^{4}$) & $1.3\times10^{4}$ ($1.2\times10^{4}$) & $1.3\times10^{5}$ ($1.3\times10^{5}$) \\
DP5 & $487$ ($1.3\times10^{3}$) & $1.8\times10^{3}$ ($7.1\times10^{3}$) & $1.3\times10^{3}$ ($2.5\times10^{3}$) & $6.6\times10^{3}$ (--) & $1.9\times10^{3}$ ($1.6\times10^{3}$) & $1.1\times10^{4}$ (--) \\
DOP853 & $476$ ($1.2\times10^{3}$) & $1.4\times10^{3}$ (--) & $1.0\times10^{3}$ ($1.5\times10^{3}$) & $2.3\times10^{3}$ ($3.5\times10^{3}$) & $958$ ($1.4\times10^{3}$) & $3.2\times10^{3}$ ($2.4\times10^{3}$) \\
GL1 & $3.3\times10^{5}$ ($1.9\times10^{6}$) & -- & $4.3\times10^{5}$ ($9.6\times10^{5}$) & -- & $5.2\times10^{5}$ ($7.4\times10^{5}$) & -- \\
GL2 & $1.4\times10^{4}$ ($4.9\times10^{4}$) & $1.1\times10^{5}$ ($3.6\times10^{5}$) & $1.5\times10^{4}$ ($2.1\times10^{4}$) & $9.7\times10^{4}$ ($1.4\times10^{5}$) & $1.4\times10^{4}$ ($1.8\times10^{4}$) & $8.9\times10^{4}$ ($1.2\times10^{5}$) \\
GL3 & $9.8\times10^{3}$ ($1.9\times10^{4}$) & $3.6\times10^{4}$ ($6.4\times10^{4}$) & $5.4\times10^{3}$ ($7.5\times10^{3}$) & $1.6\times10^{4}$ ($2.0\times10^{4}$) & $8.1\times10^{3}$ ($8.4\times10^{3}$) & $2.1\times10^{4}$ ($1.8\times10^{4}$) \\
Tao2 & $2.0\times10^{5}$ (--) & -- & $2.2\times10^{5}$ (--) & -- & $2.2\times10^{5}$ (--) & -- \\
Tao4 & $5.6\times10^{4}$ (--) & $5.6\times10^{5}$ (--) & $3.4\times10^{4}$ (--) & $3.4\times10^{5}$ (--) & $1.6\times10^{4}$ (--) & $1.6\times10^{5}$ (--) \\
\bottomrule
\end{tabular}
}
\end{table*}

The deflection angle itself is reproduced to $10^{-14}$ relative far from the hole and to
$10^{-11}$ at $b - b_c = 10^{-3}$, where the problem amplifies errors (\fig{deflection}). The
fourth-order weak-field series is good to $10^{-11}$ at $b \approx 2000$ and off by more than
1\% below $b \approx 10$; Bozza's limit is accurate near $b_c$ and fails above
$b - b_c \approx 1$. Near the critical impact parameter the number of windings follows
$-\ln(b/b_c - 1)/2\pi$, and the integration error grows exactly like the condition number
$1/(b - b_c)$: $3.3\times10^{-13}/(b - b_c)$ windings in (b) and $1.4\times10^{-10}/(b - b_c)$
in (a) (\fig{windings}). The periapsis advance agrees with $\Phi - 2\pi$ to
$4\times10^{-13}$ away from the separatrix and to $1.5\times10^{-8}$ at $p = 7.02$, where the
orbit sweeps 3.4 extra turns per period and the weak-field value $6\pi/p$ is off by 87\%
(\fig{precession}).

\begin{figure}[t]
\centering
\includegraphics[width=0.6\textwidth]{F02_deflection.pdf}
\caption{\label{fig:deflection}Deflection angle against $b - b_c$: exact, weak-field series,
Bozza's limit and the integration (DOP853, tolerance $10^{-13}$); bottom, relative errors (F2).}
\end{figure}

\begin{figure}[t]
\centering
\includegraphics[width=0.48\textwidth]{F12_windings.pdf}\hfill
\includegraphics[width=0.48\textwidth]{F13_precession.pdf}
\caption{\label{fig:windings}\label{fig:precession}Left: windings near the critical impact
parameter and their integration error, which follows the condition number $1/(b - b_c)$ (F12).
Right: periapsis advance against $p$, exact, integrated and weak-field (F13).}
\end{figure}

\subsection{Long integrations: growth of the errors (RQ1)}
\label{sec:long}

\fig{energy} shows the mass-shell error over $10^4$ radial periods of $(20, 0.5)$, with 1000
steps per period for the fixed-step methods and tolerance $10^{-10}$ for the adaptive pairs.
The hypothesis of \RQ{1} holds without exception (\tab{growth}). \textbf{RK4, DP5 and DOP853
drift linearly}: the exponents fitted over the last decade are $1.000 \pm 0.002$, and
$|\delta H|$ reaches \numdHRKfourLast{} (RK4) and \numdHDOPLast{} (DOP853).
\textbf{GL1, GL2, Tao2 and Tao4 are bounded}, with exponents below $3\times10^{-3}$; GL2 stays
at \numdHGLtwo{}. GL3's truncation error is below round-off. Its error grows from $10^{-16}$ to
\numdHGLthreeLast{} with exponent \numexpGLthree, Brouwer's random walk~\cite{brouwer1937}.

\begin{figure}[t]
\centering
\includegraphics[width=0.7\textwidth]{F07_energy_long_term.pdf}
\caption{\label{fig:energy}Mass-shell error, maximum per radial period, over $10^4$ periods
of $(20, 0.5)$ in formulation (b), drawn as its maximum over logarithmic bins (F7).}
\end{figure}

The phase error at the $N$-th periapsis (\fig{phase}) grows quadratically for the
Runge--Kutta methods (exponents 1.97--2.04) and linearly for GL1, GL2 and Tao ($1.000 \pm
0.001$). A log-log slope is a poor summary of a signed error, so \tab{growth} also fits
$\phi_N - N\Phi = c_1N + c_2N^2$. For RK4 in (b) the constant frequency error $c_1$ and the
drift $c_2$ have opposite signs and cancel near $N = \numcrossRKfour$, the dip in the figure.
GL3's phase error, \numphaseGLthreeEnd{} at $N = 10^4$, is 5 ulp of $\phi \approx
7.3\times10^4$. Its linear growth is the spacing of floating-point numbers growing with
$\phi$.

\begin{figure}[t]
\centering
\includegraphics[width=0.95\textwidth]{F09_phase_error.pdf}
\caption{\label{fig:phase}Phase error $|\phi_N - N\Phi|$ at the $N$-th periapsis in both
formulations, with slopes 1 and 2 (F9).}
\end{figure}

\begin{table*}[t]
\caption{\label{tab:growth}Growth laws over $10^4$ periods of $(20, 0.5)$ (T4): the
mass-shell error in the last period, the log-log exponent over the last decade and its
nearest law (0, 1/2, 1, 2), the same for the phase error, and the signed fit
$\phi_N - N\Phi = c_1N + c_2N^2$ (radians per orbit and per orbit squared).}
\centering\footnotesize
% Generated by scripts/make_tables.py from results/summary/. Do not edit.
\begin{tabular}{lrrlrrlrr}
\toprule
Method & $|\delta H|$, last period & Exp. & Law & Phase error & Exp. & Law & $c_1$ & $c_2$ \\
\midrule
RK4 (b) & $1.6\times10^{-8}$ & 1.00 & linear & $6.4\times10^{-5}$ & 2.04 & quadratic & $3.2\times10^{-9}$ & $-9.6\times10^{-13}$ \\
DP5 (b) & $8.4\times10^{-8}$ & 1.00 & linear & $5.0\times10^{-4}$ & 1.97 & quadratic & $-4.9\times10^{-10}$ & $-4.9\times10^{-12}$ \\
DOP853 (b) & $2.5\times10^{-7}$ & 1.00 & linear & $1.5\times10^{-3}$ & 2.00 & quadratic & $2.4\times10^{-10}$ & $1.5\times10^{-11}$ \\
GL1 (b) & $4.0\times10^{-7}$ & 0.00 & bounded & $0.93$ & 1.00 & linear & $9.3\times10^{-5}$ & $7.5\times10^{-14}$ \\
GL2 (b) & $5.0\times10^{-12}$ & 0.00 & bounded & $8.0\times10^{-6}$ & 1.00 & linear & $8.0\times10^{-10}$ & $-2.1\times10^{-20}$ \\
GL3 (b) & $6.7\times10^{-16}$ & 0.48 & random walk & $7.3\times10^{-11}$ & 1.02 & linear & $7.0\times10^{-15}$ & $2.4\times10^{-20}$ \\
Tao2 (b) & $2.6\times10^{-6}$ & 0.00 & bounded & $0.73$ & 1.00 & linear & $-7.3\times10^{-5}$ & $1.2\times10^{-12}$ \\
Tao4 (b) & $5.2\times10^{-9}$ & 0.00 & bounded & $1.8\times10^{-3}$ & 1.00 & linear & $1.8\times10^{-7}$ & $-1.0\times10^{-15}$ \\
\midrule
RK4 (a) & $8.7\times10^{-10}$ & 0.89 & linear & $8.9\times10^{-4}$ & 2.00 & quadratic & $6.7\times10^{-11}$ & $-8.9\times10^{-12}$ \\
DP5 (a) & $2.2\times10^{-7}$ & 1.00 & linear & $0.012$ & 2.01 & quadratic & $-7.1\times10^{-9}$ & $1.2\times10^{-10}$ \\
DOP853 (a) & $9.4\times10^{-8}$ & 1.00 & linear & $4.8\times10^{-3}$ & 2.00 & quadratic & $-8.4\times10^{-11}$ & $-4.8\times10^{-11}$ \\
GL1 (a) & $4.0\times10^{-7}$ & 0.00 & bounded & $0.45$ & 1.00 & linear & $-4.5\times10^{-5}$ & $8.8\times10^{-14}$ \\
GL2 (a) & $5.0\times10^{-12}$ & 0.00 & bounded & $7.5\times10^{-6}$ & 1.00 & linear & $-7.5\times10^{-10}$ & $2.2\times10^{-18}$ \\
GL3 (a) & $3.3\times10^{-16}$ & 0.25 & random walk & $2.9\times10^{-11}$ & 0.89 & linear & $-6.8\times10^{-15}$ & $2.3\times10^{-19}$ \\
\bottomrule
\end{tabular}

\end{table*}

\paragraph{Round-off.} Compensated summation decides what happens once truncation errors are
negligible (\fig{roundoff}). On the circular orbit $r_c = 10$ every Runge--Kutta map is exact,
so the error in $\phi$ after $10^4$ orbits ($10^7$ steps) is round-off alone. With compensated
summation it is \numcircComp, 7 ulp of $\phi$. With plain summation it is \numcircPlain{} and
grows like $N^{\numcircPlainExp}$: the same increment is rounded the same way at every step,
so the errors add coherently, and each is proportional to $\phi$. GL3's mass shell follows
Brouwer's law with compensation (exponent 0.48) and grows faster without it (0.67).

\begin{figure}[t]
\centering
\includegraphics[width=\textwidth]{F19_roundoff.pdf}
\caption{\label{fig:roundoff}Round-off with and without compensated summation (F19). (i) Phase
of the circular orbit $r_c = 10$, where only round-off acts. (ii) Mass shell of GL3 and RK4 on
$(20, 0.5)$. (iii) Survival on the photon sphere (Sec.~\ref{sec:unstable}).}
\end{figure}

\subsection{The role of the formulation (RQ2)}
\label{sec:formulation}

In formulation (b), $E$ and $L_z$ keep their initial values bit for bit in every sample of
every method, Tao included, for all $10^4$ periods. In (a) they are nonlinear functions of the
state, and the Runge--Kutta methods drift linearly in $E$, $L_z$ and the norm, while GL1 and
GL2 stay bounded and GL3 stays at round-off (\fig{invariants}).

\begin{figure}[t]
\centering
\includegraphics[width=\textwidth]{F08_invariants_formulation_a.pdf}
\caption{\label{fig:invariants}Formulation (a): relative errors in $E$, $L_z$ and the norm
over $10^4$ periods (F8). In (b), $E$ and $L_z$ do not change at all.}
\end{figure}

\textbf{Gauss--Legendre in (a), where it is symmetric but not symplectic, conserves as well as
in (b).} The largest mass-shell error of GL2 is $5.0138\times10^{-12}$ in both formulations,
equal in every printed digit, and the phase drifts are $-7.48\times10^{-10}$ and
$+8.03\times10^{-10}$ radians per orbit (\fig{symmetric}). For this integrable, reversible
system reversibility is enough, as the theory of reversible integrators
predicts~\cite[Ch.~XI]{gni2006}. RK4 drifts in both formulations.

\begin{figure}[t]
\centering
\includegraphics[width=0.95\textwidth]{F17_symmetric_vs_symplectic.pdf}
\caption{\label{fig:symmetric}GL2 in formulation (a), symmetric but not symplectic, against
GL2 in (b), symplectic, and RK4 in both (F17).}
\end{figure}

The formulation matters most for rays that travel far. In (a) the angular momentum
$L = r^2u^\phi$ is rebuilt from a component that falls like $r^{-2}$, which an absolute
tolerance cannot follow to full relative accuracy, so $L$ and with it $b = L/E$ drift. At
$b = 5.3$ and $r_\mathrm{far} = 10^4$, DOP853 at tolerance $10^{-13}$ makes an error of
$1.2\times10^{-8}$ in (a) and $3\times10^{-12}$ in (b); near $b_c$ the windings are 440 times
less accurate in (a) (\fig{windings}). At a fixed budget of $2\times10^4$ evaluations the
adaptive pairs reach $10^{-11}$--$10^{-12}$ in (b) for every $b$, but lose up to three orders
near $b_c$ in (a).

\paragraph{Inclined orbits.} $L^2$, the analogue of Carter's constant, is bounded for GL2 and
GL3 in both formulations and drifts linearly for RK4 and DOP853 (\fig{inclined}). The orbital
plane turns, and the turn is almost entirely a spurious precession of the node: the
inclination changes by at most $7\times10^{-9}$ for the Gauss methods, as it must with $L_z$
exact and $L^2$ bounded. The node precesses linearly in time for every method, symplectic or
not. Exact Schwarzschild has no nodal precession at all, so a drift in energy cannot change
the rate. Near the pole ($85^\circ$) every method is under-resolved at 1000 steps per period,
and (b) is the worse formulation: the node of GL2 drifts by $2.0\times10^{-5}$ radians per
period in (b) and $1.7\times10^{-7}$ in (a).

\begin{figure}[t]
\centering
\includegraphics[width=\textwidth]{F18_inclined.pdf}
\caption{\label{fig:inclined}Inclined orbits over 1000 periods (F18): (i) $L^2$ and (ii) tilt
of the orbital plane at $60^\circ$; (iii) spurious node precession against inclination. Tao4
fails at $60^\circ$ and $85^\circ$ (Sec.~\ref{sec:method-details}).}
\end{figure}

\subsection{Long integrations: cost per accuracy (RQ3)}
\label{sec:long-cost}

\fig{long-work} repeats the work-precision comparison over 1000 periods, where errors that
grow with time have time to matter. To reach a phase error of $10^{-4}$, DOP853 is still the
cheapest method (\numlongDOPfour{} evaluations per period in (b), against
\numlongGLthreeFour{} for GL3; \tab{long-cost}). \textbf{A phase error of $10^{-8}$ is reached
by GL3 alone}, for \numlongGLthreeEight{} evaluations per period. DOP853 stops at
$4.8\times10^{-6}$ at tolerance $10^{-13}$, the tightest a double-precision controller accepts,
because its phase error grows like $N^2$. For the mass shell, $|\delta H| \le 10^{-10}$,
DOP853 remains the cheapest. Over three periods (\tab{cost}) the same orbit needed
\numshortDOP{} evaluations per period from DOP853 to reach $10^{-10}$ in phase, and
\numshortGLthree{} from GL3. The ranking turns over with the length of the integration.

\begin{figure}[t]
\centering
\includegraphics[width=0.9\textwidth]{F20_long_work_precision.pdf}
\caption{\label{fig:long-work}Work-precision over 1000 radial periods of $(20, 0.5)$: phase
error and largest mass-shell error against evaluations per period (F20).}
\end{figure}

\begin{table}[t]
\caption{\label{tab:long-cost}Evaluations per radial period needed over 1000 periods of
$(20, 0.5)$ (T5b), formulation (b) with (a) in parentheses.}
\centering\small
% Generated by scripts/make_tables.py from results/summary/. Do not edit.
\begin{tabular}{lrrr}
\toprule
Method & Phase $10^{-4}$ & Phase $10^{-8}$ & $|\delta H|$ $10^{-10}$ \\
\midrule
RK4 & $6.2\times10^{3}$ ($6.6\times10^{3}$) & -- & $7.0\times10^{3}$ ($4.3\times10^{3}$) \\
DP5 & $1.5\times10^{3}$ ($3.3\times10^{3}$) & -- & $1.7\times10^{3}$ ($2.8\times10^{3}$) \\
DOP853 & $706$ ($974$) & -- & $743$ ($964$) \\
GL1 & -- & -- & -- \\
GL2 & $6.1\times10^{3}$ ($8.9\times10^{3}$) & -- & $6.1\times10^{3}$ ($7.1\times10^{3}$) \\
GL3 & $2.1\times10^{3}$ ($2.8\times10^{3}$) & $6.2\times10^{3}$ ($8.2\times10^{3}$) & $2.3\times10^{3}$ ($2.9\times10^{3}$) \\
Tao2 & -- & -- & -- \\
Tao4 & $1.5\times10^{4}$ & -- & $2.5\times10^{4}$ \\
\bottomrule
\end{tabular}

\end{table}

\subsection{Unstable and marginally stable orbits (RQ4)}
\label{sec:unstable}

\paragraph{Photon sphere.} A photon starts exactly at $r = 3$ in a plane inclined by
$45^\circ$. Its only perturbation is the integration error $\delta_0$, amplified like
$e^{\lambda_L t}$ until $|r - 3| > 0.1$ after about $\ln(0.1/\delta_0)/2\pi$ orbits. If
$\delta_0$ is proportional to the tolerance, every decade of tolerance buys
$\ln 10/2\pi = 0.37$ orbits; if $\delta_0 \propto h^p$, every decade of steps buys $0.37p$.
Both predictions hold (\fig{photon}, \tab{photon}): the adaptive pairs gain 0.33--0.43 orbits
per decade (DOP853 in (b): \numpsDOPdecade); RK4, GL2 and Tao4 gain 1.38--1.49 per decade of
steps against 1.47 predicted; GL3 gains \numpsGLthreeDecade{} against 2.20. Then GL3
saturates at 5.8--6.2 orbits whatever the step. That is the round-off limit, near the
$\operatorname{arccosh}(0.1/\epsilon_\mathrm{mach})/2\pi = 5.5$ orbits of an offset of one
machine epsilon. \textbf{No method keeps a photon on the photon sphere for more than about six
orbits in double precision.} Explicit offsets $10^{-2}$--$10^{-15}$ follow the
elliptic-integral prediction to $0.01$ orbits down to $10^{-14}$.

\begin{figure}[t]
\centering
\includegraphics[width=0.95\textwidth]{F10_photon_sphere.pdf}
\caption{\label{fig:photon}Photon sphere (F10). (a) $|r - 3|$ against coordinate time for
DOP853 at four tolerances, with the growth $e^{t/3\sqrt3}$. (b), (c) Orbits survived against
tolerance and against steps per orbit. (d) Equatorial orbits from an explicit offset against
the exact prediction.}
\end{figure}

\begin{table}[t]
\caption{\label{tab:photon}Photon-sphere survival (T6): orbits gained per decade of tolerance
or of steps per orbit, against the prediction, and the orbits survived at the finest setting.}
\centering\small
% Generated by scripts/make_tables.py from results/summary/. Do not edit.
\begin{tabular}{llrrr}
\toprule
Method & Varied & Orbits per decade & Predicted & Orbits, finest \\
\midrule
DOP853 (a) & tolerance & 0.35 & 0.37 & 4.66 \\
DOP853 (b) & tolerance & 0.33 & 0.37 & 4.59 \\
DP5 (a) & tolerance & 0.43 & 0.37 & 4.82 \\
DP5 (b) & tolerance & 0.37 & 0.37 & 4.94 \\
GL2 (a) & step size & 1.41 & 1.47 & 4.91 \\
GL2 (b) & step size & 1.44 & 1.47 & 4.75 \\
GL3 (a) & step size & 2.49 & 2.20 & 5.78 \\
GL3 (b) & step size & 2.26 & 2.20 & 5.96 \\
RK4 (a) & step size & 1.49 & 1.47 & 4.65 \\
RK4 (b) & step size & 1.45 & 1.47 & 4.56 \\
Tao4 (b) & step size & 1.38 & 1.47 & 3.84 \\
\bottomrule
\end{tabular}

\end{table}

\textbf{Plain floating-point summation makes the photon sphere look more stable than it is}
(\fig{roundoff}, panel iii). With RK4 at 8000 steps per orbit the orbit survives
\numpsPlainRKfour{} orbits with plain summation, against \numpsCompRKfour{} with compensated
summation. An increment of $r$ smaller than half an ulp of 3 is lost, $r$ stays at exactly 3,
and the radial force, which is proportional to $r - 3$, stays exactly zero: rounding freezes
the instability, the longer the smaller the step. A longer survival is not always a better
integrator.

\paragraph{ISCO.} At the ISCO the effective potential has an inflection point. With
$r = 6 + \delta$ and $L^2 = 12$, $V_\mathrm{eff} = 8/9 + \delta^3/1944 + \order(\delta^4)$, and
an integration error acts as a small constant force $\varepsilon$:
\begin{equation}
  \ddot\delta = \varepsilon - \frac{\delta^2}{1296}.
  \label{eq:isco}
\end{equation}
For $\varepsilon > 0$ a stable equilibrium appears at $\delta = 36\sqrt\varepsilon$, and the
orbit oscillates about it without plunging. For $\varepsilon < 0$ it drifts inward under a
constant acceleration until $\delta \sim -36\sqrt{|\varepsilon|}$ and then runs away in a
finite time. Both stages take a time of order $|\varepsilon|^{-1/4}$. The growth is algebraic,
and the sign of the error decides whether the orbit plunges at all. With
$\varepsilon \propto h^p$ the time to plunge should scale as $n^{p/4}$ in the steps per orbit
$n$, and as $\mathrm{tol}^{-1/4}$ for the adaptive pairs. The measurements agree
(\fig{isco}): the slope is $1.11$ for RK4 in (b), $\numiscoGLtwo$ for GL2 and
$\numiscoGLthree$ for GL3 (predicted 1, 1 and 1.5), and $\numiscoDPfive$ against tolerance for
DP5 in (b). The runs that survive oscillate on the outer side, as predicted: RK4 in (a) at
100 steps per orbit stays within $6.000 < r < 6.073$ for 2000 orbits. DOP853 never plunges at
tolerances up to $10^{-6}$; DP5 in (b) plunges at $10^{-9}$--$10^{-13}$ but not at
$10^{-6}$--$10^{-8}$.

\begin{figure}[t]
\centering
\includegraphics[width=0.95\textwidth]{F11_isco_plunge.pdf}
\caption{\label{fig:isco}Inclined ISCO orbit: orbits until $|r - 6| > 1$ against tolerance and
against steps per orbit, with the slopes fitted over the plunges (F11). Filled markers
plunged, hollow ones left outward or were still bound after 2000 orbits.}
\end{figure}

\subsection{Method-specific findings}
\label{sec:method-details}

\paragraph{Tao's coupling $\omega$.} The error of Tao4 grows in proportion to $\omega$ once
$\omega h \gtrsim 10^{-3}$, and the method diverges for $\omega h \gtrsim 1.4$; too small an
$\omega$ lets the copies separate, by up to 23 within 300 periods for $\omega = 10^{-5}$
(\fig{tao}). We use $\omega = 10^{-3}$ throughout: within a factor 1.2--5 of the best value
on the orbits tested, with the copies within $3\times10^{-6}$ over 300 periods. This choice
fails on two of the long experiments. On the ISCO the copies separate by $1.7$ in $r$ and the orbit
plunges within six orbits whatever the step; $\omega = 0.1$ holds them together. On the
inclined orbits the run is captured after 41 periods at $60^\circ$ and after 3 at $85^\circ$;
a larger $\omega$ helps at $60^\circ$ but leaves $L^2$ drifting 3000 times more than with
GL2, and no $\omega$ between $10^{-3}$ and 1 works at $85^\circ$. The coupling is a
problem-dependent parameter, and choosing it is a real cost of the method.

\paragraph{Implicit iteration.} Iterations per step grow from 2--5 at small steps to 10--23 at
the largest; GL3 needs the fewest, and (a) needs more than (b) (\fig{implicit}). Stopping the
iteration at a tolerance instead of at stagnation makes GL2 drift: over 1000 periods $|\delta
H|$ grows tenfold with tolerances $10^{-6}$ and $10^{-8}$, threefold with $10^{-10}$, and not
at all with $10^{-12}$ or at stagnation. A tolerance of $10^{-12}$ needs 3.5 iterations per
step against 6.7 for stagnation: the round-off rule is safe but not the cheapest safe rule.

\paragraph{Adaptive steps.} On the deflection the step grows roughly in proportion to $r$, and
DOP853 needs three to four times fewer steps than DP5 at the same tolerance (\fig{steps}).

\begin{figure}[t]
\centering
\includegraphics[width=0.75\textwidth]{F14_tao_omega.pdf}
\caption{\label{fig:tao}Tao4: error against the coupling $\omega$ at fixed step, and the
separation of the two copies (F14).}
\end{figure}

\begin{figure}[t]
\centering
\includegraphics[width=\textwidth]{F15_implicit.pdf}
\caption{\label{fig:implicit}Gauss--Legendre iteration (F15): iterations per step against $h$
and along the orbit, and the drift of $|\delta H|$ when the iteration stops at a tolerance.}
\end{figure}

\begin{figure}[t]
\centering
\includegraphics[width=0.75\textwidth]{F16_step_history.pdf}
\caption{\label{fig:steps}Step sizes of the adaptive pairs against $r$ on the deflection and
on an eccentric orbit (F16).}
\end{figure}
